% 图契约：恒流大圆环在抗磁平面上方的悬浮；图中的端口、几何、箭头和方程对应本题。
% 版面：沿用前三模块的几何图+推导分区；目标印宽17cm，简图至中等密度。
% 布局：左侧真实结构，右侧关键关系；真实电路边不作为架构流程箭头处理。
\documentclass[tikz,border=7pt]{standalone}
\usepackage{fontspec,xeCJK,amsmath,amssymb}
\IfFontExistsTF{Times New Roman}{\setmainfont{Times New Roman}}{\setmainfont{TeX Gyre Termes}}
\IfFontExistsTF{Heiti SC}{\setCJKmainfont[AutoFakeBold=1.5]{Heiti SC}}{\setCJKmainfont[AutoFakeBold=1.5]{FandolHei-Regular}}
\usetikzlibrary{arrows.meta,bending,calc,patterns,positioning,decorations.pathmorphing}
\definecolor{ink}{HTML}{243745}
\definecolor{blue}{HTML}{275DAB}
\definecolor{green}{HTML}{177C63}
\definecolor{red}{HTML}{BD4935}
\definecolor{orange}{HTML}{AA6C16}
\definecolor{purple}{HTML}{715393}
\tikzset{every node/.style={font=\fontsize{9.3}{12.5}\selectfont,text=ink},
arrow/.style={-{Stealth[length=5pt 1.5,width'=0pt 3]},shorten >=2pt,shorten <=1pt,line width=1pt},
arrow short/.style={-{Stealth[length=3pt 1.0,width'=0pt 2.5]},shorten >=1pt,shorten <=0pt,line width=.8pt},
axis/.style={arrow,draw=ink!60,line width=.6pt},
vec/.style={arrow,draw=blue},
wire/.style={draw=ink!75,line width=.8pt},
guide/.style={dashed,draw=ink!45,line width=.5pt},
card/.style={rounded corners=2pt,draw=blue!65,fill=blue!4,align=center,text width=6.5cm,inner xsep=9pt,inner ysep=7pt},
title/.style={font=\fontsize{10.5}{14}\selectfont\bfseries,text width=15cm,align=center},
label/.style={font=\fontsize{8.4}{11}\selectfont,align=center},
note/.style={label,text width=6.1cm}}
\begin{document}
\begin{tikzpicture}[x=1cm,y=1cm]
\path[use as bounding box] (-8,-3.5) rectangle(8,3.65);
\node[title] at (0,3.25) {M5-164\quad 恒流大圆环在抗磁平面上方的悬浮};
\begin{scope}[shift={(-4.05,0)}]
\fill[blue!4](-2.8,-2.4)rectangle(2.8,0);\draw[ink,thick](-2.8,0)--(2.8,0);\node[label]at(-1.6,-.3){完全抗磁平面};\draw[red,thick](0,1.2)ellipse(1.4 and .25);\draw[purple,dashed](0,-1.2)ellipse(1.4 and .25);\draw[vec,green](0,1.2)--(0,2.2)node[right]{$\mathbf m$};\draw[vec,purple](0,-1.2)--(0,-2.2)node[right]{$\mathbf m_{\rm im}$};\draw[<->,orange](2,0)--(2,1.2)node[midway,right]{$h$};\node[above]at(-1.4,1.5){实物};\node[below]at(-1.4,-1.5){镜像};\end{scope}
\node[card,draw=blue!65,fill=blue!4] at (3.9,1.8) {$r\gg h,\quad F_z\simeq\mu_0I^2r/(2h)$};
\node[card,draw=green!65,fill=green!4] at (3.9,0) {$h_0=\mu_0I^2r/(2Mg)$};
\node[card,draw=orange!65,fill=orange!4] at (3.9,-1.8) {$T=2\pi\sqrt{h_0/g}$};
\end{tikzpicture}\end{document}